\chapter{Adapting the Workflow without BSDRunner}

The Rails, PostgreSQL, nginx, and Minitest boundaries in this book do not
depend on a graphical interface. The provisioning exercises do. BSDRunner
adds guarded workflows around Blackship and the host firewall; removing it
means accepting responsibility for those workflows rather than merely changing
the spelling of a few commands.

This appendix is a translation guide, not a parallel tutorial. A reader using
raw Blackship should first create an equivalent private network and five jail
definitions, then resume the application exercises only after the addresses,
outbound policies, and dependency graph match the book.

\section{Action-to-concept mapping}

\begin{description}
  \item[Create a Service Jail] Write and validate a jail definition, build its
    root, and run \texttt{blackship up NAME}. BSDRunner's starter also installs
    packages, applies service-specific settings, and guards generated
    credentials.
  \item[Manifest editor] Edit the appropriate per-jail TOML fragment. The
    shared manifest owns project storage and network definitions; a fragment
    owns one jail's desired state.
  \item[Dry Run] Use \texttt{blackship eva --dry-run NAME} to classify the
    difference between the installed definition and a running jail before
    applying a hot update or restart.
  \item[Apply] Validate and install the candidate TOML, then use
    \texttt{blackship eva NAME} for supported live changes. Restart-required
    changes must go through the dependency-aware restart path.
  \item[Start, stop, restart] Use \texttt{blackship up},
    \texttt{blackship down}, and \texttt{blackship restart}. Blackship expands
    dependencies and dependents; do not substitute unrelated \texttt{jail(8)}
    commands and expect the same graph behavior.
  \item[Console and logs] Use \texttt{blackship console NAME} and
    \texttt{blackship logs NAME}. The Rails rc.d services also keep their own
    application logs inside the catalog and store jails.
  \item[Health] Use \texttt{blackship health NAME} for a one-shot result or
    add \texttt{--watch} for continuous monitoring.
  \item[Port forwarding] Use \texttt{blackship expose} only after deciding
    which host interface and port should be reachable. BSDRunner additionally
    renders, validates, installs, and rolls back its managed PF profile.
\end{description}

\section{Loading the split manifest}

Blackship accepts repeated \texttt{--config} arguments. A POSIX shell can
assemble the shared manifest and every jail fragment without relying on shell
arrays:

\begin{lstlisting}[language=sh]
set -- --config /usr/local/etc/bsdrunner/blackship.toml
for fragment in /usr/local/etc/bsdrunner/blackship.d/*.toml; do
  [ -f "$fragment" ] || continue
  set -- "$@" --config "$fragment"
done

blackship "$@" check
blackship "$@" eva --dry-run catalog
blackship "$@" health catalog
\end{lstlisting}

The paths retain BSDRunner's project layout because that is the layout built
in this book. A standalone Blackship project may use different paths, but the
ordering rule is the same: load the shared configuration first and the
per-jail fragments afterward.

\section{What must be replaced deliberately}

Without BSDRunner, provide your own equivalent for typed destructive-action
confirmation, atomic manifest installation, PF validation and rollback,
one-time credential handling, address-collision checks, and readable operation
logs. Those protections are part of the tutorial's operating assumptions, not
cosmetic UI details.

Once an equivalent jail graph is healthy, Chapters 5 through 14 apply directly
to the two Rails applications, Chapters 16 through 20 apply to the web and
service boundaries, and Chapters 22 through 26 apply to the external test and
its automation contract.
